\section{Fixed-Dimensional Risk Bounds for Exact Orthogonal Actions}
\label{app:spectral-risk}

This appendix separates the best risk in an entire action class from the error of selected fitted operators. The class and metric below are narrower than arbitrary linear prediction: they impose exact group laws and orthogonality after population second-moment whitening. The result is a mathematical extension, not an additional empirical result or a certificate for the recorded FULL/AFS scores.

\subsection{Observation contract and invariant-subspace budget}
\label{sp:setup}
Let a finite group $G$, of order $m$, act on a finite state set $S$, of size $p$; the action may be nonfree or nonfaithful. A context $c$ supplies the complete labeled block $\{z_{c,s}:s\in S\}$ in fixed coordinates $z_{c,s}\in\R^{1\times r}$, with $r\ge1$. Expectations use $C\sim P_C$ and an independent uniform $s\in S$; group averages include every element, including the identity. Assume $\Sigma=\mathbb E[z_{C,s}^{\top}z_{C,s}]\succ0$, and put $x=z\Sigma^{-1/2}$, so $\mathbb E[x^{\top}x]=\Id_r$. This is a second moment; zero mean is not required. The group product $gh$ means first $h$, then $g$, giving the row-action convention
\[
 \mathcal R_r(G)=\{\rho:G\to O(r):\rho_e=\Id_r,\ \rho_g\rho_h=\rho_{hg}\}.
\]
Define two population risks, with $x_g(c,s)=x(c,g\cdot s)$:
\begin{equation}
 \eta=\inf_{\{A_g\}}\frac1{rm}\sum_g\mathbb E\norm{xA_g-x_g}_2^2,
 \qquad L^\star=\inf_{\rho\in\mathcal R_r(G)}L(\rho),\quad
 L(\rho)=\frac1{rm}\sum_g\mathbb E\norm{x\rho_g-x_g}_2^2.
 \label{sp:risks}
\end{equation}
The first class allows independent unconstrained $A_g\in\R^{r\times r}$; the second requires one exact $r$-dimensional orthogonal action shared across contexts. In raw coordinates $R_g=\Sigma^{-1/2}\rho_g\Sigma^{1/2}$, its loss is
\[
 \frac1{rm}\sum_g\mathbb E\norm{(zR_g-z_g)\Sigma^{-1/2}}_2^2.
\]
This is neither the row-normalized metric $D$ in \cref{eq:D} nor the later training-energy normalization. Exact zero-error actions preserve $\Sigma$ and become orthogonal after whitening, but positive-error optima over all $GL(r)$ need not do so. Accordingly, a lower bound below excludes only the class in \cref{sp:risks}. Independently fitted Ridge and finite-step penalized Joint generally do not belong to that exact class. Their recorded errors are neither covered by its lower bound nor upper-bound witnesses for $L^\star$.

On the real Hilbert space $\mathcal H=L^2(P_C\times\mathrm{Unif}(S))$, define
\[
 (Ja)(c,s)=x_{c,s}a,\qquad (U_gf)(c,s)=f(c,g\cdot s),\qquad
 P=JJ^*,\qquad T=\frac1m\sum_gU_gPU_g^*.
\]
Then $J^*J=\Id_r$, $P$ is a rank-$r$ orthogonal projection, and $U_gU_h=U_{hg}$ with each $U_g$ orthogonal. The group average $T$ commutes with every $U_g$, satisfies $0\preceq T\preceq\Id$ and $\operatorname{tr}T=r$, and is supported on
\[
 \mathcal W=\operatorname{span}\{U_gJa:g\in G,\ a\in\R^r\},\qquad \dim\mathcal W\le mr.
\]
Define
\begin{equation}
 C_r(T)=\max_{Q}\operatorname{tr}(TQ),\qquad \gamma_r=C_r(T)/r,
 \quad
 \begin{gathered}
 Q^2=Q=Q^*,\quad [Q,U_g]=0\ \text{for all }g,\\[-2pt]
 \operatorname{rank}Q\le r,\quad\operatorname{range}Q\subseteq\mathcal W.
 \end{gathered}
 \label{sp:budget}
\end{equation}
Thus $0\le\gamma_r\le1$. Since invariant $Q$ satisfies $\operatorname{tr}(PQ)=\operatorname{tr}(TQ)$,
\[
 r-C_r(T)=\min_Q\norm{(\Id-Q)J}_{\mathrm{HS}}^2.
\]
This is a best invariant-subspace approximation quantity, not a new notion of identifiability. Spectral constructions for invariant approximation have precedents \citep{barbieri2020approximation}; the connection to the exact-action risk is established below.

The rank constraint must respect \emph{real} irreducible types. If type $\tau$ has real dimension $d_\tau$ and multiplicity $m_\tau$, the corresponding self-adjoint multiplicity operator has eigenvalues $\lambda_{\tau,1}\ge\cdots\ge\lambda_{\tau,m_\tau}\ge0$, each repeated $d_\tau$ times in the real space. With integer multiplicities $k_\tau$, standard real representation theory \citep{serre1977linear} gives
\begin{equation}
 C_r(T)=\max_{\substack{0\le k_\tau\le m_\tau\\\sum_\tau d_\tau k_\tau\le r}}
 \sum_\tau d_\tau\sum_{j=1}^{k_\tau}\lambda_{\tau,j}.
 \label{sp:integer}
\end{equation}
Indeed, the commutant on each multiplicity space is a real, complex or quaternionic matrix algebra. Its self-adjoint operator diagonalizes, and the variational principle selects the largest $k_\tau$ eigenvalues for a fixed multiplicity budget. Optimizing the resulting integer real-dimension budget proves \cref{sp:integer}. A nontrivial real $C_7$ rotation plane cannot be split into two separately admissible one-dimensional directions.

\begin{spectrallemma}[Unconstrained regression risk]
\label{sp:regression}
The optimal unconstrained matrices are $A_g=J^*U_gJ$, and
\begin{equation}
 \eta=\frac1{rm}\sum_g\norm{(\Id-P)U_gJ}_{\mathrm{HS}}^2
 =1-\operatorname{tr}(T^2)/r.
 \label{sp:eta}
\end{equation}
\end{spectrallemma}
\begin{proof}
Since $J$ is an isometry, $JA_g=PU_gJ$ is the least-squares projection of $U_gJ$ onto $\operatorname{range}J$. Averaging its residual squared gives $r-\operatorname{tr}(PT)$. Group averaging is the Hilbert--Schmidt orthogonal projection onto the commutant, so $\operatorname{tr}(PT)=\operatorname{tr}(T^2)$. Division by $r$ proves the result.
\end{proof}

\subsection{The fixed-dimensional risk bound}
\begin{spectraltheorem}[Risk of the entire exact orthogonal action class]
\label{sp:main}
Under the assumptions of \cref{sp:setup},
\begin{equation}
 \boxed{\quad 2(1-\gamma_r)\ \le\ L^\star\ \le\ 2(1-\gamma_r^2).\quad}
 \label{sp:sandwich}
\end{equation}
The lower bound holds for every $\rho\in\mathcal R_r(G)$, and some member of this same $r$-dimensional class satisfies the upper bound.
\end{spectraltheorem}
\begin{proof}
For a fixed $\rho$, define on maps $A:\R^r\to\mathcal H$
\[
 \Pi_\rho(A)=\frac1m\sum_gU_gA\rho_g^*.
\]
Each summand is a Hilbert--Schmidt isometry. Reindexing by inverses and products proves that the average is self-adjoint and idempotent; it projects onto the intertwiner space $\{A:U_gA=A\rho_g\ \forall g\}$. With $B=\Pi_\rho(J)$, expansion gives
\begin{equation}
 L(\rho)=2-\frac2r\langle J,B\rangle_{\mathrm{HS}}
        =2-\frac2r\norm B_{\mathrm{HS}}^2.
 \label{sp:projection-risk}
\end{equation}
Let $Q_B$ project onto $\operatorname{range}B$. This range is invariant, contained in $\mathcal W$, and has rank at most $r$. Cauchy--Schwarz yields
\[
 \norm B_{\mathrm{HS}}^2=\langle Q_BJ,B\rangle_{\mathrm{HS}}
 \le\norm{Q_BJ}_{\mathrm{HS}}\norm B_{\mathrm{HS}},\qquad
 \norm B_{\mathrm{HS}}^2\le\operatorname{tr}(PQ_B)
 =\operatorname{tr}(TQ_B)\le C_r(T).
\]
The zero-$B$ case is immediate. Substitution into \cref{sp:projection-risk} proves the lower bound.

For the upper bound, choose a maximizer $Q$ in \cref{sp:budget}, of rank $k\le r$. In an auxiliary space put
\[
 \widetilde{\mathcal H}=\mathcal H\oplus\R^{r-k},\quad
 \widetilde U_g=U_g\oplus\Id_{r-k},\quad
 \widetilde J=(J,0),\quad \widetilde Q=Q\oplus\Id_{r-k}.
\]
Let $V:\R^r\to\widetilde{\mathcal H}$ be an orthonormal basis of $\operatorname{range}\widetilde Q$. An SVD of $V^*\widetilde J$ and orthogonal Procrustes alignment give another such basis $J'$ with
\[
 \langle J',\widetilde J\rangle_{\mathrm{HS}}
 =\norm{V^*\widetilde J}_*
 \ge\norm{V^*\widetilde J}_F^2
 =\operatorname{tr}(\widetilde Q\widetilde J\widetilde J^*)=C_r(T).
\]
The inequality uses singular values in $[0,1]$, since $V$ and $\widetilde J$ are isometries. Invariance of $\operatorname{range}J'$ makes $\rho_g=J'^*\widetilde U_gJ'$ an exact orthogonal action on $\R^r$ and makes $J'$ an intertwiner. Hence, for $B'=\Pi_\rho(\widetilde J)$,
\[
 C_r(T)\le\langle J',\widetilde J\rangle_{\mathrm{HS}}
 =\langle J',B'\rangle_{\mathrm{HS}}
 \le\sqrt r\norm{B'}_{\mathrm{HS}}.
\]
Equation~\eqref{sp:projection-risk} gives $L(\rho)\le2-2C_r(T)^2/r^2$. The data have zero auxiliary coordinates, so this construction changes neither the original loss nor the predictor dimension; the final matrices remain $r\times r$.
\end{proof}

\paragraph{Relation to existing theory.}
\label{sp:positioning}
The averaging, intertwiner and dimension-accounting tools are standard in approximate representation theory \citep{gowers2017inverse,moore2015approximate}. The statement here connects invariant approximation to \emph{fixed-dimensional prediction risk}; it does not claim priority for those tools. Sun's realizability and error decomposition concern compatible transformations through an encoder \citep{sun2026transformationlaws}; EbC identifies jointly learned latent coordinates and actions under generative and diversity assumptions \citep{schmidt2025equivariance}. Neither observation contract is replaced by this frozen-coordinate bound.

\subsection{Finite observable form and an insufficient-summary example}
\label{sp:observable}
Let $D=mr$, $v(c,s)=m^{-1/2}(x_{c,g\cdot s})_{g\in G}\in\R^{1\times D}$ and $K=\mathbb E[v^\top v]$. Its blocks are $K_{g,h}=m^{-1}\mathbb E[x_g^\top x_h]$. For $a=(a_g)_g$, define
\[
 Wa=m^{-1/2}\sum_gU_gJa_g,\qquad (\Lambda_ta)_k=a_{kt^{-1}}.
\]
Direct multiplication gives $WW^*=T$, $W^*W=K$, and $U_tW=W\Lambda_t$. Thus $K$ commutes with the known block action $\Lambda$, has trace $r$, and is a positive contraction. The polar decomposition of $W$ identifies $\operatorname{range}K$ with $\mathcal W$ equivariantly: their nonzero spectra agree \emph{with their real irreducible types}. Therefore the optimization in \cref{sp:budget}, now using $K,\Lambda$ on $\R^D$, gives
\[
 C_r(K)=C_r(T),\qquad \eta=1-\operatorname{tr}(K^2)/r.
\]
Zero-eigenvalue blocks add no trace objective. This construction uses complete same-context labeled orbits as a structural diagnostic; its information is not the input available to a single-activation predictor. Computing the statistic still requires the type-constrained optimization, not ordinary top-$r$ eigenvalue truncation.

\paragraph{Registration and the full action average.}
A group automorphism merely reindexes the terms of the full uniform group average: the image action, $T$ and its invariant-subspace budget are unchanged. Conjugating the state registration to a different image subgroup can change the budget, but need not do so. For $C_7$ there are $6!=720$ seven-cycle permutations, generating 120 distinct cyclic subgroups; these are different counting units. Any comparison of their budgets is initially descriptive. Removing identity or registration-dependent fixed points changes the averaging contract and cannot be substituted into \cref{sp:main}.

\paragraph{Energy by type and its contribution to the budget.}
\label{sp:type-energy}
Let $\Pi_\tau$ project onto the real $\tau$-isotypic subspace of the known action $\Lambda$, and let $Q^\star$ attain $C_r(K)$. Distinguish total energy from energy retained by the dimension budget:
\[
 e_\tau=\frac{\operatorname{tr}(\Pi_\tau K)}r,\qquad
 b_\tau=\frac{\operatorname{tr}(\Pi_\tau KQ^\star)}r,\qquad
 \sum_\tau e_\tau=1,\quad\sum_\tau b_\tau=\gamma_r.
\]
Here $0\le b_\tau\le e_\tau$; selected shares can depend on the choice of a tied maximizer even though $\gamma_r$ does not. The trivial type describes components invariant under the state action, including possible contextual and global-constant components, rather than uniquely identifying context information. A large $\gamma_r$ can also coexist with stable state coding when trivial energy is small. Removing trivial or weak directions changes the declared coordinates, dimension or metric; reporting their contributions preserves the target.

\paragraph{Equal regression risk and eigenvalues can conceal different action risk.}
\label{sp:counterexample}
Take $G=S=C_2\times C_2\times C_7$ acting by addition, uniform states $(a,b,k)$, one context, and $r=1$. Set
\[
 x_A(a,b,k)=\frac{(-1)^a+(-1)^b}{\sqrt2},\qquad
 x_B(a,b,k)=\sqrt2\cos(2\pi k/7).
\]
Both have zero mean and unit second moment. For $g=(u,v,w)$ their optimal scalar regression coefficients are
\[
 a_A(g)=\tfrac12\big((-1)^u+(-1)^v\big),\qquad
 a_B(g)=\cos(2\pi w/7).
\]
Both satisfy $\mathbb E_g a(g)^2=1/2$, hence $\eta=1/2$. For $A$, $T$ has eigenvalue $1/2$ on each of two distinct real one-dimensional characters; for $B$, it is $\Id_2/2$ on a real irreducible rotation plane. The nonzero eigenvalues are thus the same, $\{1/2,1/2\}$. An exact real one-dimensional action is a character $\chi:G\to\{\pm1\}$ and is necessarily trivial on $C_7$. Its loss is $2-2\mathbb E_g[\chi(g)a(g)]$. Character orthogonality makes the largest correlation $1/2$ for $A$ and zero for $B$, so $L_A^\star=1$ and $L_B^\star=2$; correspondingly $\gamma_1=1/2$ and $0$. This illustrates the standard representation-dimension obstruction: regression risk and an untyped spectrum are insufficient summaries. The complete labeled observations distinguish the examples, so this is not observational nonidentifiability of the full data.

\paragraph{Scale of the type-budget correction for C7 and D4.}
\label{sp:type-scale}
For these two groups all real irreducible types have dimension one or two. Define the ordinary spectral budget $C_r^{\mathrm{untyped}}(K)$ as the sum of the largest $r$ eigenvalues of $K$, including zeros if needed. Order tied eigenvectors by complete irreducible blocks. This top-$r$ selection can split at most one two-dimensional block at its boundary. Dropping the selected half yields a feasible invariant subspace of dimension at most $r$ and loses at most one eigenvalue, which is at most one because $K\preceq\Id$. Hence
\[
 0\le \Delta_{\mathrm{type}}
 :=C_r^{\mathrm{untyped}}(K)-C_r(K)\le1,
 \qquad \Delta_{\mathrm{type}}/r\le1/r.
\]
This elementary consequence gives a normalized upper bound of $1/16$ at $r=16$; it is not an observed effect or a claim that the bound is tight. A zero type correction would leave the class-risk bound valid. The one-dimensional separation above therefore does not predict a large correction on the LLM panel, and does not justify searching dimensions for a positive correction.

For the four-corner action $R(i)=i+1$, $S(i)=1-i\pmod4$, the state permutation representation decomposes as $\chi_{+,+}\oplus\chi_{-,-}\oplus E$, where $\chi_{a,b}(R)=a$, $\chi_{a,b}(S)=b$ and $E$ is the real two-dimensional type. Indeed, the stabilizer of corner zero is $\{e,R^3S\}$, so a one-dimensional type has multiplicity $(1+ab)/2$. Context copies do not introduce the absent types: $K$ has zero energy in $\chi_{+,-}$ and $\chi_{-,+}$. This is an analytical input check for any future numerical diagnostic, not a completed numerical check. The corner action has four states even though the group has eight elements.

\subsection{A finite-context certificate and its practical limitation}
\label{sp:certificate}
Suppose $C_1,\ldots,C_n$ are independent identically distributed complete context blocks, independent of the choice of coordinates and preprocessing. Assume the \emph{population} bound $\norm{z_{C,s}}_2\le H$ almost surely for every state, with $H$ fixed in advance. Define
\[
 \Sigma_i=p^{-1}\sum_sz_{C_i,s}^\top z_{C_i,s},\quad
 v_z(c,s)=m^{-1/2}(z_{c,g\cdot s})_g,\quad
 M_i=p^{-1}\sum_sv_z(C_i,s)^\top v_z(C_i,s),
\]
and sample means $\widehat\Sigma,\widehat M$ with population means $\Sigma,M$. When $\widehat\Sigma\succ0$, let
\[
 \widehat K=(\Id_m\otimes\widehat\Sigma^{-1/2})\widehat M
                    (\Id_m\otimes\widehat\Sigma^{-1/2}).
\]
Complete uniform state averaging preserves the fixed $\Lambda$ symmetry exactly and gives $\operatorname{tr}\widehat K=r$.

\begin{spectraltheorem}[Conservative complete-block certificate]
\label{sp:finite}
For $\alpha\in(0,1)$, set
\[
 \tau_n=H^2\sqrt{\frac8n\log\frac{2(D+r)}\alpha},\qquad
 a_n=\lambda_{\min}(\widehat\Sigma)-\tau_n.
\]
If $a_n>0$, define $\varepsilon_n=\tau_n/a_n+H^2\tau_n/a_n^2$, $\widehat\gamma_r=C_r(\widehat K)/r$, and
\[
 \gamma_- =\max\{0,\widehat\gamma_r-\varepsilon_n\},\qquad
 \gamma_+ =\min\{1,\widehat\gamma_r+\varepsilon_n\}.
\]
There is an event of probability at least $1-\alpha$ on which, whenever this precheck passes,
\begin{equation}
 2(1-\gamma_+)\le L^\star\le2(1-\gamma_-^2).
 \label{sp:confidence}
\end{equation}
If $a_n\le0$, the procedure returns inconclusive; discarding weak directions changes the specified $r$-dimensional target. This bounds the unconditional probability of issuing an incorrect interval, not the error probability conditional on passing the precheck.
\end{spectraltheorem}
\begin{proof}
Each $\Sigma_i,M_i$ lies between zero and $H^2\Id$, so its centered square is bounded by $H^4\Id$. Applying the two-sided matrix Hoeffding inequality \citep[Theorem~1.3]{tropp2012user} and a union bound for dimensions $r,D$ gives, with probability at least $1-\alpha$,
\[
 \norm{\widehat\Sigma-\Sigma}_{\rm op}\le\tau_n,\qquad
 \norm{\widehat M-M}_{\rm op}\le\tau_n.
\]
On this event and when $a_n>0$, both second moments are at least $a_n\Id$. For positive $A,B\succeq a\Id$, the inverse-square-root resolvent formula gives
\[
 \norm{A^{-1/2}-B^{-1/2}}_{\rm op}
 \le\frac{\norm{A-B}_{\rm op}}\pi\int_0^\infty\frac{t^{-1/2}}{(a+t)^2}\,dt
 =\frac{\norm{A-B}_{\rm op}}{2a^{3/2}}.
\]
Writing $A_0=\Id_m\otimes\Sigma^{-1/2}$ and $\widehat A_0=\Id_m\otimes\widehat\Sigma^{-1/2}$, expand
\[
 K-\widehat K=A_0(M-\widehat M)A_0
 +(A_0-\widehat A_0)\widehat M A_0
 +\widehat A_0\widehat M(A_0-\widehat A_0).
\]
Using $\norm{\widehat M}_{\rm op}\le H^2$ yields $\norm{K-\widehat K}_{\rm op}\le\varepsilon_n$. Every feasible invariant projection has rank at most $r$, so
\[
 |\operatorname{tr}((K-\widehat K)Q)|\le r\varepsilon_n,
 \qquad |C_r(K)-C_r(\widehat K)|\le r\varepsilon_n.
\]
The feasible set is fixed by $G,\Lambda,r$, not selected from the data. Apply \cref{sp:main} to obtain \cref{sp:confidence}.
\end{proof}

\paragraph{Why this does not certify the present experiments.}
The bound is formally valid but can be uninformative. Necessarily
\[
 \lambda_{\min}(\widehat\Sigma)\le\operatorname{tr}(\widehat\Sigma)/r\le H^2/r,
 \qquad a_n>0\ \Longrightarrow\
 n>8r^2\log\frac{2r(m+1)}\alpha.
\]
At $r=16$ and $\alpha=0.05$, even the possibility of passing the precheck requires at least 17,492 independent complete contexts for $C_7$, or 17,734 for $D_4$. These are necessary conditions for this conservative formula, not sufficient conditions, intrinsic sample-complexity lower bounds or proposed experimental sizes. Even after passing, \cref{sp:confidence} may be $[0,2]$.

The current 40-template panel and its overlapping partitions cannot meet this condition. States, actions, layers and repeated splits are not additional independent contexts. A sample maximum cannot replace the stipulated population bound $H$; clipping would change the target. Data-dependent preprocessing requires independent fitting data or a separate analysis. Multiple reported model/layer targets require an explicit simultaneous-error allocation; selecting the best layer is not covered by the single-target formula. No population certificate is reported here.

\paragraph{Finite-bank interpretation.}
\label{sp:finite-bank}
A fixed, complete coordinate bank can instead define $P_C$ as its declared empirical distribution. If its second moment is positive definite, \cref{sp:observable,sp:main} then give deterministic structural summaries and risk bounds for that bank, without an independent-context assumption. This uses the second moment in the fixed input coordinates, not an additional centering step; all states and actions, including identity and fixed points, receive their declared uniform weights. Computing the whitening metric on this bank defines a new descriptive target, not a held-out score for the historical predictors. Its bounds are not confidence intervals for unseen contexts, and numerical rank or conditioning failures cannot be repaired by silently changing the target dimension. No finite-bank spectral values are reported as empirical results in this paper.

The accompanying C7/D4 implementation computes the budget and type energies, with separate checks for input provenance, numerical whitening, projector completeness, real-type multiplicities and agreement with direct least squares. Its analytic software fixtures are implementation checks, not statistical calibration or LLM evidence. No exact-action witness is constructed. A future witness would need to realize the proof's invariant subspace and satisfy the group laws; the stated upper bound applies to a correctly constructed witness on the same bank and metric. Scoring a frozen predictor on different contexts or with historical $D$ would be a separate prediction comparison.
